% Hướng dẫn cấu hình hệ thống ERP Sao Việt Nhật: một file chung cho mọi module.
% Biên dịch: xelatex cau-hinh-he-thong.tex (chạy hai lần cho mục lục).
% Tài liệu sử dụng từng phòng link sang các mốc \label{...} trong file này.
\documentclass[11pt,a4paper]{article}
\usepackage{fontspec}
\usepackage[vietnamese]{babel}
\setmainfont{Times New Roman}
\newfontfamily\tieude{Arial}
\usepackage[a4paper,margin=1.9cm,top=1.7cm]{geometry}
\usepackage{graphicx}
\usepackage{enumitem}
\usepackage{xcolor}
\usepackage{titlesec}
\usepackage{tocloft}
\usepackage{needspace}
\usepackage{fancyhdr}
\usepackage[hidelinks]{hyperref}
\usepackage[most]{tcolorbox}
\usepackage{longtable}
\usepackage{booktabs}
\graphicspath{{hinh/}}

% Màu cam gạch của phần mềm.
\definecolor{gach}{HTML}{C2410C}
\definecolor{xamnen}{HTML}{F4F4F5}
\definecolor{xamchu}{HTML}{52525B}

\setlist{nosep,leftmargin=1.5em}
\setlength{\parskip}{0.4em}
\setlength{\parindent}{0pt}
\setlength{\emergencystretch}{2em}

% Chương = \section, mục = \subsection: tài liệu chạy liền, không ngắt trang ở mỗi chương.
\titleformat{\section}{\tieude\Large\color{gach}}{Chương \thesection.}{0.4em}{}[{\color{gach}\titlerule[0.8pt]}]
\titlespacing*{\section}{0pt}{1.2em}{0.6em}
\titleformat{\subsection}{\tieude\large}{\thesubsection}{0.6em}{}
\titlespacing*{\subsection}{0pt}{1.1em}{0.3em}
\renewcommand{\cftsecpresnum}{Chương }
\setlength{\cftsecnumwidth}{5.2em}
\renewcommand{\cftsecfont}{\tieude}
\renewcommand{\cftsecpagefont}{\tieude}
\setlength{\cftbeforesecskip}{0.5em}
\addto\captionsvietnamese{\renewcommand{\contentsname}{Mục lục}}
\renewcommand{\cfttoctitlefont}{\tieude\large}
\setlength{\cftbeforetoctitleskip}{0pt}
\setlength{\cftaftertoctitleskip}{0pt}

\pagestyle{fancy}
\fancyhf{}
\fancyhead[L]{\tieude\footnotesize\color{xamchu}Hướng dẫn cấu hình hệ thống}
\fancyhead[R]{\tieude\footnotesize\color{xamchu}ERP Sao Việt Nhật}
\fancyfoot[C]{\tieude\footnotesize\thepage}
\renewcommand{\headrulewidth}{0.4pt}

% Đoạn bối cảnh đầu mỗi chương/mục: việc này là gì, ai làm, khi nào làm.
\newtcolorbox{boicanh}{colback=xamnen,colframe=xamnen,boxrule=0pt,arc=3pt,
  left=9pt,right=9pt,top=6pt,bottom=6pt,before skip=4pt,after skip=8pt}

% Tên ô, tên nút trên màn hình.
\definecolor{vang}{HTML}{FEF08A}
\newcommand{\oo}[1]{{\setlength{\fboxsep}{1.2pt}\colorbox{vang}{\strut #1}}}
% Đường đi trong thanh bên.
\newcommand{\duong}[1]{{\tieude\small #1}}
% Số khoanh đỏ trên ảnh.
\newcommand{\so}[1]{{\tieude\bfseries\small\textcolor{white}{\colorbox{red!80!black}{\kern-2pt#1\kern-2pt}}}}
% Ảnh có viền mảnh.
\newcommand{\khung}[2]{{\setlength{\fboxsep}{0pt}\setlength{\fboxrule}{0.5pt}%
  \fcolorbox{black!20}{white}{\includegraphics[width=\dimexpr#1-1pt\relax]{#2}}}}
% Nháy kép thẳng.
\newcommand{\q}[1]{\textquotedbl#1\textquotedbl{}}
% Ảnh bên trái, chú thích bên phải: \canh{tỉ lệ rộng ảnh}{file}{nội dung}.
\newcommand{\canh}[3]{%
  \par\noindent
  \begin{minipage}[t]{#1\linewidth}\vspace{0pt}\khung{\linewidth}{#2}\end{minipage}\hfill
  \begin{minipage}[t]{\dimexpr0.97\linewidth-#1\linewidth\relax}\vspace{0pt}\setlength{\parskip}{0.5em}#3\end{minipage}\par\medskip}
% Ảnh rộng hết dòng.
\newcommand{\rong}[2][\linewidth]{\par\noindent\begin{center}\khung{#1}{#2}\end{center}\vspace{-0.6em}}

\begin{document}

% ---------- Đầu tài liệu ----------
\thispagestyle{empty}
\begin{tcolorbox}[colback=gach,colframe=gach,arc=4pt,left=14pt,right=14pt,top=14pt,bottom=12pt]
\color{white}\tieude
{\footnotesize ERP SAO VIỆT NHẬT}\par\vspace{4pt}
{\fontsize{24}{28}\selectfont Hướng dẫn cấu hình hệ thống}\par\vspace{6pt}
{\small Phòng ban, phân quyền, tài khoản, danh mục dùng chung và danh mục của từng phòng}\par\vspace{8pt}
{\footnotesize Cập nhật 31/10/2026}
\end{tcolorbox}

\begin{boicanh}
Tài liệu này gom mọi việc khai báo ban đầu của phần mềm vào một chỗ. Người làm là quản trị phần mềm và người
được giao khai danh mục ở mỗi phòng. Dùng khi xưởng mới bắt đầu chạy phần mềm, và mỗi khi có phòng, người, máy,
giấy hay công đoạn mới. Tài liệu sử dụng của từng phòng sẽ trỏ về đúng mục ở đây thay vì viết lại.

Ảnh trong tài liệu chụp từ phần mềm thật. Ô đỏ đánh số là chỗ cần bấm hoặc cần điền, số khớp với lời giải
thích bên cạnh.
\end{boicanh}

\tableofcontents

% ============================================================
\section{Thứ tự khai}\label{thu-tu-khai}

\begin{boicanh}
Nhiều danh mục cần danh mục khác có trước. Máy phải thuộc một tổ, công đoạn phải chọn máy chạy được, phiếu
tính giá phải chọn khách và giấy. Khai ngược thứ tự thì ô chọn sẽ trống. Sơ đồ dưới cho biết khai cái gì
trước, cái gì sau.
\end{boicanh}

\rong[0.84\linewidth]{luong-cau-hinh.png}

\begin{enumerate}
\item Phòng ban và tổ sản xuất. Vai trò thuộc phòng, máy thuộc tổ, nên phòng ban khai đầu tiên.
\item Vai trò và ma trận quyền của từng phòng (chương \ref{phan-quyen}).
\item Chấm công (điểm chấm công, ca làm, lịch và ngày lễ), loại nghỉ phép, cấu hình lương. Hồ sơ nhân viên gán
ca từ đây, bảng lương đầu tiên tính theo luật ở đây.
\item Hồ sơ nhân sự: tạo tài khoản, gán vai, khai mức lương từng người. Lập nhóm dùng chung nếu cần.
\item Danh mục nền: Đơn vị \& quy đổi, rồi Thiết bị \& Máy móc, rồi Giấy và Vật tư khác, cuối cùng là Công đoạn.
\item Khách hàng, rồi lập thử một phiếu tính giá mẫu so với giá tính tay.
\end{enumerate}

Thành phẩm và khuôn không cần khai trước. Thành phẩm tự sinh khi chốt đơn hàng bán. Phiếu tính giá chỉ hỏi
khuôn có sẵn hay làm mới; khuôn bế chọn ở bước lệnh sản xuất. Khai báo kho, Tiêu chí KCS, Xe giao hàng, Mục thu chi
khai trước khi bắt đầu dùng kho, KCS, giao hàng, thu chi (mục \ref{danh-muc-khuon} trở đi).

% ============================================================
\section{Phân quyền}\label{phan-quyen}

\begin{boicanh}
Phân quyền quyết định mỗi người mở được màn nào, sửa được gì và thấy dữ liệu của ai. Chương này giải thích
cách đọc và cách gán quyền, dùng chung cho mọi module. Quyền riêng của từng màn nằm ngay đầu mục khai màn đó
ở các chương sau. Người làm là quản trị phần mềm, hoặc trưởng phòng được bật quyền sửa ma trận.
\end{boicanh}

Quyền gắn vào \emph{vai trò}, mỗi vai thuộc một phòng. Một người dùng được màn nào phải đủ ba thứ: phòng có vai
đó, vai đó bật quyền trên ma trận, và tài khoản của người đó đã được gán vai.

\subsection{Mở ma trận quyền}\label{phan-quyen-ma-tran}

\begin{boicanh}
Mỗi phòng tự có danh sách vai của mình. Ví dụ Phòng Kinh Doanh có ba vai: Trợ lý kinh doanh, sale và Trưởng
phòng kinh doanh. Làm mục này khi phòng mới lập, khi cần thêm vai, hoặc khi muốn mở thêm hay thu bớt quyền
của một vai. Sửa ma trận là sửa cho mọi người đang giữ vai đó.
\end{boicanh}

\duong{Nhân sự \& Lương > Phòng ban}, chọn phòng ở cây bên trái. Ảnh là Phòng Kinh Doanh.

\rong[0.74\linewidth]{20-vai-tro.png}

\so{1} Mở tab \oo{Vai trò \& Quyền}.

\so{2} Bấm chip tên vai để mở ma trận của vai đó. Số trên chip là số người đang giữ vai.

\so{3} \oo{+ Vai trò} tạo vai mới trong phòng, gõ tên vào ô \oo{Tên vai trò} rồi bật quyền.

\so{4} Sửa xong bấm \oo{Lưu thay đổi}. Chưa bấm thì chưa có hiệu lực.

Cần vai gần giống vai có sẵn thì bấm \oo{Nhân bản vai trò} rồi sửa chỗ khác nhau, đỡ phải bật lại từng dòng.

\Needspace*{12\baselineskip}
\subsection{Đọc một dòng quyền}\label{phan-quyen-dong}

\begin{boicanh}
Mỗi dòng của ma trận là một màn trên thanh bên, dòng nào cũng có cùng một bộ ô. Đọc hiểu một dòng là đọc được
cả ma trận. Ví dụ dưới là dòng Tính giá của vai Trợ lý kinh doanh.
\end{boicanh}

\rong{21-dong-quyen.png}

\begin{enumerate}[leftmargin=2.2em,labelsep=0.5em,itemsep=2pt]
\item[\so{1}] Ô tick: bật hay tắt cả màn với vai này.
\item[\so{2}] \oo{Xem}: mở màn và xem dữ liệu. Tắt Xem thì Thao tác cũng tắt theo.
\item[\so{3}] \oo{Thao tác}: thêm, sửa, xoá. Bật Thao tác thì Xem tự bật theo.
\item[\so{4}] \oo{Phạm vi}: thấy dữ liệu của ai, giải thích ở mục \ref{phan-quyen-pham-vi}.
\item[\so{5}] Chip \oo{N/M chi tiết}: bấm để mở quyền chi tiết. Màn không có quyền chi tiết thì không có chip.
\item[\so{6}] Quyền chi tiết. Ở Tính giá là \oo{Xem chi tiết giá vốn}: tắt thì chỉ thấy giá vốn tổng.
\end{enumerate}

Sáu màn luôn bật Xem, không tắt được: Nội quy, Hồ sơ của tôi, Chấm công, Nghỉ phép, Tăng ca, Lương.

\subsection{Phạm vi}\label{phan-quyen-pham-vi}

\begin{boicanh}
Phạm vi cho biết người đó thấy dữ liệu của ai: \oo{Của tôi} là của mình, \oo{Cả phòng} là của cả phòng mình
(gồm tổ con), \oo{Tất cả} là cả công ty.
\end{boicanh}

Ở khối Kinh doanh, khách và phiếu tính giá, báo giá, đơn của khách đó đều tính theo \oo{NV phụ trách} của
khách. Cấu hình đang dùng ở Phòng Kinh Doanh:

\begin{itemize}
\item Sale để \oo{Của tôi}: chỉ thấy khách mình phụ trách.
\item Trợ lý kinh doanh để \oo{Cả phòng}: thấy mọi khách của các sale trong phòng.
\end{itemize}

Khách nào cũng phải có \oo{NV phụ trách}. Khách bỏ trống ô này thì chỉ người để \oo{Tất cả} mới thấy.

\subsection{Nhóm dùng chung}\label{phan-quyen-nhom}

Chỉ cần khi muốn một người thấy khách của vài sale chứ không phải cả phòng. Người để \oo{Của tôi} thấy thêm
khách của người cùng nhóm. Trợ lý đã để \oo{Cả phòng} thì không cần nhóm.

\Needspace*{22\baselineskip}
\subsection{Gán vai và lập nhóm}\label{phan-quyen-gan-vai}

\begin{boicanh}
Có vai rồi thì phải gán vai cho người. Tab Nhân sự của phòng cho làm nhiều người một lần, nên dùng khi mới
khai phòng hoặc khi đổi vai cho cả nhóm người. Nhóm dùng chung cũng lập ở đây.
\end{boicanh}

\duong{Nhân sự \& Lương > Phòng ban}, chọn phòng.

\rong[0.64\linewidth]{22-gan-vai-nhom.png}

\so{1} Mở tab \oo{Nhân sự}.

\so{2} Tick những người cần làm. Tick bao nhiêu người cũng được, thanh trên cùng báo số đã chọn.

\so{3} \oo{Gán vai trò}: chọn một vai của phòng cho những người đã tick.

\so{4} \oo{Thêm vào nhóm}: mở danh sách nhóm ở ảnh dưới.

\so{5} Chip tím dưới tên người là nhóm họ đang ở. Badge xanh bên phải là vai.

Vai và nhóm gắn theo tài khoản. Người chưa có tài khoản sẽ bị bỏ qua khi gán.

\canh{0.42}{23-them-vao-nhom.png}{%
Sau khi bấm \oo{Thêm vào nhóm}:

\so{1} Danh sách nhóm có sẵn. Tick nhóm để thêm những người đã chọn vào, bấm lại nhóm đã tick để gỡ họ ra.
Nhóm nào mới có một phần người đã chọn thì ghi \q{1/2 người đã ở}.

\so{2} \oo{Tạo nhóm mới từ N người này} khi chưa có nhóm phù hợp.}

\Needspace*{22\baselineskip}
\subsection{Tài khoản của một người}\label{phan-quyen-tai-khoan}

\begin{boicanh}
Mỗi hồ sơ nhân sự có một tài khoản đăng nhập. Mục này dùng khi chỉ cần xem hay đổi vai của một người, hoặc
khi người đó quên mật khẩu, nghi lộ mật khẩu, cần chặn đăng nhập tạm thời.
\end{boicanh}

\duong{Nhân sự \& Lương > Hồ sơ nhân sự}, mở hồ sơ người đó.

\rong[0.8\linewidth]{24-tai-khoan.png}

\so{1} Mở tab \oo{Tài khoản \& Quyền}.

\so{2} Ô \oo{Vai trò} chỉ liệt kê vai của phòng ghi trên hồ sơ.

Khối \oo{Bảo mật} có ba nút: \oo{Đặt lại mật khẩu}, \oo{Thu hồi mọi phiên} và \oo{Khóa tài khoản}. Thu hồi
mọi phiên là đăng xuất người đó trên mọi máy. Người đã nghỉ việc tự không đăng nhập được, không cần khoá tay.

% ============================================================
\Needspace*{14\baselineskip}
\section{Nhân sự}\label{nhan-su}

\begin{boicanh}
Chương này dựng bộ khung người của công ty: phòng ban và tổ sản xuất, rồi hồ sơ từng nhân viên kèm tài khoản
đăng nhập. Người làm là phòng Hành chính nhân sự hoặc quản trị phần mềm. Vai trò và nhóm dùng chung đã nói ở
chương \ref{phan-quyen}.
\end{boicanh}

\subsection{Quyền của hai màn}\label{nhan-su-quyen}

Mở ma trận như mục \ref{phan-quyen-ma-tran}. Dòng \oo{Phòng ban}: \oo{Thao tác} là thêm, sửa, xoá phòng và vai
trò. Bấm chip chi tiết có thêm ba quyền:

\rong[0.86\linewidth]{30-quyen-phong-ban.png}

\begin{itemize}
\item \oo{Đặt trưởng phòng}: chọn hoặc đổi người đứng đầu của phòng, tổ.
\item \oo{Đổi cấp trên (cây tổ chức)}: chuyển một phòng, tổ sang trực thuộc phòng khác. Lúc tạo phòng mới thì
chọn cấp trên không cần quyền này.
\item \oo{Sửa ma trận phân quyền}: sửa và lưu ma trận, nhân bản vai, lập nhóm dùng chung. Không bật thì chỉ
xem được ma trận.
\end{itemize}

Hai quyền đầu phải bật kèm \oo{Thao tác}. Quyền thứ ba dùng riêng được.

\Needspace*{16\baselineskip}
Dòng \oo{Hồ sơ nhân sự} có mười quyền chi tiết. Phạm vi của màn này chỉ có \oo{Cả phòng} và \oo{Tất cả}.

\rong[0.86\linewidth]{31-quyen-nhan-su.png}

\begin{itemize}
\item \oo{Xem lương \& BHXH}: không bật thì ô lương, bảo hiểm, thuế trên hồ sơ bị ẩn.
\item \oo{Sửa lương \& BHXH}: bật ô này thì \oo{Xem lương \& BHXH} tự bật theo. Tắt xem lương thì ô này tắt theo.
\item \oo{Đổi trạng thái nhân viên}: lên chính thức, tạm nghỉ, đình chỉ, nghỉ việc.
\item \oo{Điều chuyển \& đổi chức danh}: chuyển người sang phòng, tổ khác.
\item \oo{Duyệt yêu cầu cập nhật}: duyệt đề nghị sửa thông tin nhân viên gửi từ \oo{Hồ sơ của tôi}.
\item \oo{Xuất Excel danh sách}: cột lương chỉ có khi được xem lương.
\item \oo{Đặt lại mật khẩu}, \oo{Khóa / Mở tài khoản}, \oo{Thu hồi phiên}, \oo{Gán vai trò}: bốn việc trên tài
khoản đăng nhập, xem mục \ref{phan-quyen-tai-khoan}.
\end{itemize}

\subsection{Tạo phòng, tổ}\label{nhan-su-tao-phong}

\begin{boicanh}
Phòng ban là gốc của mọi thứ: vai trò thuộc phòng, máy thuộc tổ, lương khoán tính theo tổ. Tổ sản xuất cũng là
một phòng, chỉ khác là nằm dưới phòng khác trên cây.
\end{boicanh}

\canh{0.5}{32-them-phong.png}{%
\duong{Nhân sự \& Lương > Phòng ban}. Phòng cấp cao nhất thì bấm \oo{Thêm phòng gốc}. Tổ nằm trong một phòng
thì rê chuột lên phòng đó ở cây bên trái, bấm dấu \oo{+}.

\so{1} Gõ tên phòng hoặc tổ.

\so{2} \oo{Đơn vị trực thuộc}: phòng cha. Để trống là phòng gốc. Bấm dấu \oo{+} ở cây thì ô này tự điền.

\so{3} Bấm \oo{Tạo phòng ban}.}

\subsection{Trưởng phòng và khối}\label{nhan-su-sua-phong}

\begin{boicanh}
Tạo xong, mở lại phòng để đặt người đứng đầu và đánh dấu phòng thuộc khối nào. Khối quyết định phòng đó hiện
ở phân hệ nào: người khối Kinh doanh mới được chọn làm NV phụ trách khách, phòng khối Sản xuất mới lên phân hệ
Sản xuất.
\end{boicanh}

\canh{0.5}{33-sua-phong.png}{%
Chọn phòng ở cây, bấm \oo{Chỉnh sửa}.

\so{1} \oo{Đơn vị Trực thuộc}: đổi phòng cha. Cần quyền \oo{Đổi cấp trên}.

\so{2} \oo{Người đứng đầu}: chỉ chọn được người thuộc phòng này hoặc tổ con của nó.

\so{3} \oo{Bộ phận thuộc Khối Sản xuất}: cả cây con tự tính là sản xuất.

\so{4} \oo{Bộ phận thuộc Khối Kinh doanh}: cả cây con tự tính là kinh doanh. Người trong khối hiện ở ô
\oo{NV phụ trách} của màn Khách hàng.

\so{5} \oo{Bộ phận Giao hàng}: chỉ đúng phòng, tổ này, không lan xuống tổ con. Người trong đó là tài xế, hiện ở
tab Nhân viên giao hàng và được tính khoán km.

\so{6} \oo{Tổ in}: chỉ đúng tổ này. Thợ in đi làm Chủ nhật, lễ thì ăn công thay vì tính sản lượng.

Xong bấm \oo{Lưu thay đổi} ở chân hộp.

Ô \oo{Tỷ lệ Lương thử việc} ở dưới chỉ để xem, số này sửa ở Cấu hình lương (mục \ref{nhan-su-cau-hinh-luong}).}

\Needspace*{18\baselineskip}
\subsection{Chấm công: điểm, ca, ngày lễ}\label{nhan-su-cham-cong}

\begin{boicanh}
Chấm công cần ba thứ có trước: chỗ được bấm chấm công, ca làm của từng người, và lịch tuần làm việc kèm ngày
lễ. Thiếu ca thì phần mềm không biết ai đi muộn, thiếu lịch thì công chuẩn của tháng sai. Khai cả ba trước khi
thêm nhân viên, vì hồ sơ nhân viên gán ca nền từ đây.
\end{boicanh}

Dòng \oo{Chấm công} trong ma trận: \oo{Xem} và \oo{Thao tác} luôn bật cho mọi người (tự bấm giờ, xem công của
mình). Việc cấu hình nằm ở ba quyền chi tiết:

\rong[0.86\linewidth]{3d-quyen-cham-cong.png}

\begin{itemize}
\item \oo{Điểm chấm công}, \oo{Khai ca}, \oo{Lịch \& Ngày lễ}: mỗi quyền mở đúng một tab cùng tên. Muốn thêm,
sửa thì phạm vi của dòng phải là \oo{Tất cả}.
\item \oo{Bảng công tháng}, \oo{Xem Nhật ký chấm công}: xem công và từng lượt chấm của người khác.
\item Các quyền duyệt phiếu đi muộn, về sớm, nghỉ nửa buổi, \oo{Duyệt yêu cầu chỉnh công}, \oo{Chấm bù / sửa công},
\oo{Chốt kỳ công / Mở lại kỳ}: việc hằng tháng của người quản lý công, không phải cấu hình.
\end{itemize}

\canh{0.42}{3d-diem-them.png}{%
\duong{Nhân sự \& Lương > Chấm công}, tab \oo{Điểm chấm công}, bấm \oo{Thêm điểm chấm công}.

\so{1} \oo{Tên điểm chấm công}, ví dụ Xưởng in.

\so{2} Đứng tại xưởng thì bấm \oo{Lấy vị trí GPS hiện tại của tôi}, vĩ độ và kinh độ tự điền.

\so{3} Hoặc bấm lên bản đồ để đặt tâm.

\so{4} \oo{Bán kính cho phép (mét)}: nhân viên chỉ bấm được \oo{Chấm công} khi đứng trong vòng này.

\so{5} \oo{Lưu cấu hình}.

Công ty có nhiều xưởng thì khai mỗi xưởng một điểm.}

\canh{0.42}{3d-ca-form.png}{%
Tab \oo{Khai ca}, mục \oo{A · Ca làm việc}, bấm \oo{Thêm ca làm việc}.

\so{1} \oo{Tên ca làm việc}, ví dụ Ca hành chính.

\so{2} \oo{Dung sai đi muộn}: vào trễ trong số phút này chưa tính đi muộn.

\so{3} \so{4} \oo{Giờ vào ca}, \oo{Giờ ra ca}. Ca có nghỉ trưa thì bật \oo{Có nghỉ giữa ca}. Ca qua nửa đêm thì
bật \oo{Ca qua đêm}.

\so{5} \oo{Phụ cấp cơm}, tính theo ngày.

\so{6} \oo{Phụ cấp ca}, tính theo ca.

\so{7} \oo{Lưu ca làm việc}.

Giờ làm thực của ca phải bằng giờ công chuẩn (mục \ref{nhan-su-cau-hinh-luong}), lệch thì phần mềm báo đỏ và
không cho lưu.}

\Needspace*{20\baselineskip}
Mục \oo{B · Phân ca tháng} bên dưới gán ca cho từng người, từng ngày:

\rong[0.9\linewidth]{3d-phan-ca.png}

\so{1} \oo{Đặt ca chính cho tất cả}: gán một ca nền cho mọi người trong lưới. Ô để trống thì lấy ca nền.
\so{2} \oo{Bút ca}: chọn một ca (hoặc \oo{Nghỉ}) rồi bấm, kéo trên lưới để tô. Phím 1 đến 9 chọn ca, phím 0 là
nghỉ. \so{3} Bấm \oo{Lưu}.

\Needspace*{20\baselineskip}
\rong[0.9\linewidth]{3f-ca-nen-luoi.png}

Người chưa có ca nền thì chưa chấm công được. Thêm nhân viên xong phải đặt ca nền cho họ. \so{1} Chip
\oo{Chưa có ca} lọc ra những người còn thiếu, số trong ngoặc là số người. \so{2} Bấm ô ở cột \oo{Ca nền} của một
người để đặt riêng cho người đó:

\canh{0.5}{3f-ca-nen-form.png}{%
Chọn \oo{Ca}, chọn ngày \oo{Áp dụng từ}, bấm \oo{Áp dụng}. Ca nền áp từ ngày đó về sau cho mọi tháng, khác với tô
ca trên lưới chỉ đúng ngày đã tô. Chỉ đặt được từ hôm nay trở đi, vì đổi ca ngày đã chấm công sẽ làm sai lương
đã tính.

Chọn \oo{bỏ gán ca nền} để gỡ ca nền của người đó.}

\canh{0.55}{3d-lich.png}{%
Tab \oo{Lịch \& Ngày lễ}:

\so{1} \oo{Tuần làm việc chuẩn}: bật các thứ đi làm. Thứ tắt là ngày nghỉ tuần.

\so{2} \oo{Lưu tuần làm việc}.

\so{3} \oo{Thêm ngày}: khai ngày lễ, ngày làm bù (ảnh dưới).

Công chuẩn của tháng phần mềm tự tính: ngày làm trong tuần, trừ ngày lễ, cộng ngày làm bù. Phần \oo{Xem trước
tháng} cho thấy kết quả.}

\canh{0.42}{3d-ngay-le.png}{%
\so{1} \oo{Ngày}.

\so{2} \oo{Tên ngày}, ví dụ Quốc khánh.

\so{3} \oo{Loại ngày}: \oo{Nghỉ lễ} (được nghỉ), \oo{Làm bù} (đi làm vào ngày lẽ ra được nghỉ), \oo{Nghỉ không
lương}.

\so{4} \oo{Hưởng nguyên lương}: bật thì mọi người được cộng một công ngày đó.

Bấm \oo{Thêm ngày} để lưu.}

\Needspace*{18\baselineskip}
\subsection{Loại nghỉ phép}\label{nhan-su-nghi-phep}

\begin{boicanh}
Nhân viên xin nghỉ phải chọn một loại nghỉ. Loại nghỉ quyết định ngày nghỉ có được trả lương không và mỗi năm
được nghỉ tối đa bao nhiêu ngày.
\end{boicanh}

\rong[0.86\linewidth]{3d-quyen-nghi-phep.png}

Dòng \oo{Nghỉ phép}: \oo{Quản danh mục loại nghỉ} mới được thêm, sửa loại nghỉ. \oo{Duyệt đơn nghỉ phép} để
duyệt đơn người khác. \oo{Huỷ đơn nghỉ của mình} nên bật cho mọi nhân viên.

\canh{0.45}{3d-loai-nghi.png}{%
\duong{Nhân sự \& Lương > Nghỉ phép}, tab \oo{Loại nghỉ}, bấm \oo{Thêm loại nghỉ mới}.

\so{1} Tên, ví dụ Phép năm, Nghỉ ốm.

\so{2} \oo{Hạn mức / năm (số ngày)}. Nhập 0 nếu không giới hạn, ví dụ nghỉ không lương.

\so{3} \oo{Hưởng nguyên lương}: đơn được duyệt thì ngày nghỉ tự chấm công ký hiệu P.

\so{4} \oo{Đang sử dụng}: tắt thì nhân viên không chọn được loại này khi xin nghỉ.

\so{5} \oo{Lưu cấu hình}.}

\Needspace*{18\baselineskip}
\subsection{Cấu hình lương}\label{nhan-su-cau-hinh-luong}

\begin{boicanh}
Cấu hình lương là luật chung để phần mềm tính bảng lương: giờ công chuẩn, hệ số tăng ca, các khoản thu nhập,
bảo hiểm, thuế. Khai một lần, tháng nào cũng dùng. Mức lương của từng người khai riêng ở mục
\ref{nhan-su-luong-nv}.
\end{boicanh}

\rong[0.86\linewidth]{3d-quyen-luong.png}

Dòng \oo{Lương}: \oo{Xem cấu hình lương} để xem, bật thêm \oo{Thao tác} để sửa. \oo{Lương nhân viên} mở tab
mức lương từng người, muốn sửa cũng phải kèm \oo{Thao tác}. Các quyền còn lại (\oo{Bảng lương tháng}, \oo{Chốt
bảng lương / Mở lại kỳ}, \oo{Đánh dấu đã chi lương}, \oo{Duyệt tạm ứng}, \oo{Xuất bảng lương / file chuyển
khoản}) là việc hằng tháng.

\Needspace*{20\baselineskip}
\duong{Nhân sự \& Lương > Lương}, tab \oo{Cấu hình lương}. Tab con \oo{Cơ chế lương theo bộ phận}:

\rong[0.68\linewidth]{3d-chl-chung.png}

\so{1} Tab con \oo{Cơ chế lương theo bộ phận}. \so{2} \oo{Giờ công chuẩn / ngày}: lương chia cho số giờ này ra
giá một giờ công. Ô \oo{Ca phải khớp giờ công chuẩn} bên dưới bắt mọi ca khai ở tab Khai ca phải đúng số giờ
này. \so{3} \oo{\% lương thử việc}: áp cho mọi người đang thử việc. Luật lao động yêu cầu tối thiểu 85\%, phần
mềm vẫn cho lưu số thấp hơn nhưng báo đỏ. \so{4} \oo{Công tối thiểu để hưởng cơm / phụ cấp ca}: ngày làm ít
hơn số công này thì không được phụ cấp cơm, phụ cấp ca.

Hai ô không khoanh: \oo{Hạn mức chỉnh công / tháng} là số ngày mỗi người được tự xin chỉnh công trong một tháng
(đếm theo ngày, không theo số đơn; HCNS chấm bù trực tiếp thì không bị giới hạn; 0 là không giới hạn).
\oo{Công tối thiểu để tạm ứng / nhận lương đợt 1}: từ đầu kỳ tới ngày lập phiếu phải có ít nhất ngần này công
(đi làm, phép có lương, lễ) mới lập được phiếu tạm ứng hoặc lương đợt 1, kể cả phiếu nhân viên tự xin. 0 là tắt
điều kiện.

\Needspace*{20\baselineskip}
Kéo xuống:

\rong[0.68\linewidth]{3d-chl-bo-phan.png}

\so{1} \so{2} Hệ số tăng ca: ngày thường, ngày nghỉ tuần, ngày lễ (mặc định 150\%, 200\%, 300\%), trả theo giờ
trên đơn giá giờ của mức nền. Các ô còn lại trong khối:
\begin{itemize}[nosep,leftmargin=1.2em]
\item \oo{Trần giờ tăng ca / tháng}: hết trần thì không tạo được phiếu tăng ca nữa, phiếu đang chờ duyệt cũng
tính. 0 là tắt trần. \oo{Trần giờ một phiếu tăng ca}: độ dài tối đa một phiếu.
\item \oo{Làm nguyên công - ngày nghỉ tuần}, \oo{Làm nguyên công - ngày lễ}: đi làm trọn ngày vào ngày nghỉ
tuần, ngày lễ thì được trả theo hệ số này thay cho tăng ca.
\item \oo{Tiền một suất cơm tăng ca} (0 là tắt) và \oo{Tăng ca bao nhiêu phút thì được một suất cơm}. Cơm tăng
ca tính riêng, không đụng cơm ca.
\item \oo{Phụ cấp làm ban đêm}: cộng thêm cho giờ làm từ 22h tới 6h. \oo{Phụ cấp tăng ca đêm}: cộng thêm cho
giờ tăng ca rơi vào 22h tới 6h.
\end{itemize}
\so{3} \oo{Đổi bộ phận}: chọn một phòng, tổ để khai phần riêng của nó.
\so{4} \oo{Chuyên cần}: bật cho cả tổ thì người trong tổ được tính thưởng chuyên cần. Số tiền khai theo từng
người ở mục \ref{nhan-su-luong-nv}.

\Needspace*{20\baselineskip}
Tab con \oo{Danh mục khoản thu nhập}:

\rong[0.68\linewidth]{3d-khoan-thu-nhap.png}

\so{1} \oo{Thêm khoản}: tạo các khoản như trợ cấp nhà ở, phụ cấp xăng xe, trang phục. \so{2} \oo{Chịu thuế}:
tick nếu khoản này tính vào thu nhập chịu thuế thu nhập cá nhân. Tạo khoản ở đây mới chỉ là bước 1. Bước 2 là
gán khoản và số tiền cho từng người ở tab \oo{Lương nhân viên}.

\Needspace*{20\baselineskip}
Tab con \oo{Bảo hiểm \& Thuế}:

\rong[0.68\linewidth]{3d-bao-hiem.png}

\so{1} Tỷ lệ người lao động đóng BHXH, BHYT, BHTN, trần lương đóng bảo hiểm, đoàn phí công đoàn, số ngày nghỉ
không lương thì thôi đóng BHXH trong tháng. Đổi khi luật đổi. Thêm hai ô: \oo{TNLĐ-BNN (công ty đóng)} chỉ dùng
cho người đã có nơi khác đóng bảo hiểm, công ty chỉ chịu khoản này, không trừ vào lương. \oo{Trần khấu trừ kỷ
luật}: tiền phạt, bồi thường trừ vào lương không quá ngần này phần trăm lương thực trả; 0 là trừ trọn số đã ghi.

\Needspace*{20\baselineskip}
Cuối tab là thuế thu nhập cá nhân:

\rong[0.68\linewidth]{3d-thue.png}

\so{1} \oo{Giảm trừ bản thân}. \so{2} \oo{Giảm trừ mỗi người phụ thuộc}. \so{3} Biểu thuế lũy tiến từng bậc.
Người có cách tính thuế \oo{Khấu trừ 10\%} (hợp đồng dưới 3 tháng, thời vụ, thực tập) không theo biểu này: thu
nhập chịu thuế cả tháng từ \oo{Ngưỡng bắt đầu khấu trừ tại nguồn} trở lên thì bị trừ theo \oo{Thuế suất khấu trừ
tại nguồn}, không được giảm trừ gia cảnh. Cách tính thuế của từng người xem ở mục \ref{nhan-su-luong-nv}.

\Needspace*{16\baselineskip}
\rong[0.68\linewidth]{3f-khau-tru-tre.png}

Cuối tab là bảng \oo{khấu trừ đi trễ / về sớm}: mỗi lần đi trễ hay về sớm không phép, quá dung sai của ca, bảng
lương tự trừ một khoản. Mỗi bậc là \q{tới N phút thì trừ X đồng một lần}, dòng cuối để trống phút là áp cho mọi
số phút lớn hơn. Bấm \oo{Thêm bậc} để thêm dòng.

Tab con \oo{Khoán km giao hàng} khai mức tiền theo km cho xe giao hàng, chỉ cần khi công ty có đội xe.

\subsection{Thêm nhân viên}\label{nhan-su-them}

\begin{boicanh}
Mỗi người làm ở công ty có một hồ sơ, kể cả công nhân xưởng. Hồ sơ là nơi tính chấm công, nghỉ phép, lương.
Tài khoản đăng nhập là tuỳ chọn: công nhân xưởng có thể không cần.
\end{boicanh}

\duong{Nhân sự \& Lương > Hồ sơ nhân sự}, bấm \oo{Thêm nhân viên}. Hộp thêm có năm bước.

\canh{0.56}{34-them-nv.png}{%
\so{1} Năm bước. Bấm \oo{Tiếp} để sang bước sau, bước nào chưa cần thì bỏ qua.

\so{2} \oo{Họ tên}, bắt buộc.

\so{3} \oo{Phòng/Tổ}, bắt buộc. Vai trò ở bước 5 lấy theo phòng này.

\so{4} \oo{Chức danh}.

\so{5} \oo{Trạng thái}: để \oo{Thử việc} thì phải điền \oo{Ngày hết thử việc}. Tới ngày đó phần mềm nhắc
chuyển chính thức.

\so{6} \oo{Tiếp}.}

\canh{0.56}{35-buoc-tai-khoan.png}{%
Bước 5 \oo{Tài khoản}:

\so{1} Tick ô \oo{Tạo tài khoản đăng nhập}. Không tick thì chỉ tạo hồ sơ.

\so{2} \oo{Tên đăng nhập}.

\so{3} \oo{Mật khẩu tạm}: đưa cho nhân viên để đăng nhập lần đầu.

\so{4} \oo{Vai trò}: để trống thì người đó đăng nhập được nhưng chưa thấy màn nào.

\so{5} \oo{Lưu \& xem hồ sơ}.}

\subsection{Cấp tài khoản cho hồ sơ đã có}\label{nhan-su-cap-tai-khoan}

\canh{0.56}{36-cap-tai-khoan.png}{%
Hồ sơ tạo trước mà chưa có tài khoản thì mở hồ sơ đó.

\so{1} Tab \oo{Tài khoản \& Quyền}. Phần mềm báo \oo{Chưa có tài khoản đăng nhập}.

Gõ \oo{Tên đăng nhập}. Ô \oo{Mật khẩu ban đầu} phần mềm tự điền sẵn, muốn đổi thì bấm \oo{Tạo mật khẩu
khác}. Chọn \oo{Vai trò}.

\so{2} Bấm \oo{Cấp tài khoản}.}

\Needspace*{18\baselineskip}
\subsection{Mức lương từng người}\label{nhan-su-luong-nv}

\begin{boicanh}
Người chưa khai mức lương thì bảng lương tính 0 đồng cho người đó. Khai ngay sau khi thêm nhân viên, và khai lại
khi tăng lương.
\end{boicanh}

\rong[0.9\linewidth]{3d-lnv.png}

\duong{Nhân sự \& Lương > Lương}, tab \oo{Lương nhân viên}. \so{1} Chip \oo{Chưa khai lương} lọc ra những người
còn thiếu. \so{2} Bấm \oo{Thiết lập} ở dòng của người đó.

\canh{0.5}{3d-lnv-form.png}{%
Tab \oo{Mức lương \& Phụ cấp}:

\so{1} \oo{Lương cơ bản}, bắt buộc.

\so{2} \oo{Lương trách nhiệm}. Mức nền bằng cơ bản cộng trách nhiệm.

\so{3} \oo{Thưởng chuyên cần}: để 0 nếu không có. Nghỉ ngày nào trừ dần theo luật ở Cấu hình lương.

\so{4} \oo{Cách trả lương}: \oo{Lương khoán} so lương theo công với tiền khoán sản lượng rồi lấy bên cao hơn.
\oo{Công nhật} trả theo công và giờ tăng ca, không nhận tiền khoán.

Dưới cùng là các khoản thu nhập đã tạo ở danh mục, điền số tiền cho người này.

\so{5} \oo{Lưu điều chỉnh}.}

\Needspace*{18\baselineskip}
\rong[0.7\linewidth]{3f-sua-luong-bh.png}

Cùng hộp đó, tab \oo{Bảo hiểm \& Công đoàn}: \so{1} \oo{Mức đóng BHXH}, bắt buộc, là số ghi trên hợp đồng bảo
hiểm của người này. \so{2} \oo{Bảo hiểm đóng ở nơi khác}: tick khi người này đã có nơi khác đóng BHXH, BHYT,
BHTN; công ty thôi trừ ba khoản đó, chỉ đóng TNLĐ-BNN. \so{3} \oo{Đoàn viên công đoàn}: chỉ người tick mới bị trừ
đoàn phí, bằng tỷ lệ ở Cấu hình lương nhân mức nền (cơ bản cộng trách nhiệm); người thử việc không đóng.
\so{4} \oo{Áp dụng giảm trừ bản thân khi tính thuế TNCN}: bỏ tick nếu người này đã đăng ký giảm trừ bản thân ở
nơi làm việc khác, vì luật chỉ cho đăng ký ở một nơi. Phía dưới hộp có bảng mô phỏng số tiền người lao động và
công ty đóng mỗi tháng để soát lại.

\Needspace*{18\baselineskip}
\rong[0.8\linewidth]{3f-ho-so-luong.png}

Thông tin để tính thuế và trả lương nằm ở hồ sơ: \duong{Hồ sơ nhân sự}, mở người đó, tab \oo{Lương \& BHXH}, bấm
\so{1} \oo{Sửa lương \& BHXH} (bấm xong nút đổi thành \oo{Hủy sửa}). \so{2} \oo{Người phụ thuộc}: mỗi người được
giảm trừ thêm khi tính thuế, mức ở Cấu hình lương. \so{3} \oo{Cách tính thuế TNCN}: mặc định \oo{Luỹ tiến từng
phần}. Người hợp đồng dưới 3 tháng, thời vụ, thực tập chọn \oo{Khấu trừ 10\% tại nguồn}; người đã nộp mẫu cam kết
08 chọn \oo{Có cam kết 08/CK-TNCN}. \so{4} \oo{Số tài khoản} và \so{5} \oo{Ngân hàng}: thiếu thì dòng của người
này trong file chuyển khoản lương bị trống. Hai ô \oo{Số sổ BHXH}, \oo{MST cá nhân} điền luôn nếu có. Bấm
\oo{Lưu}. Nút \oo{Sửa lương \& BHXH} chỉ hiện cho người có quyền sửa lương của Hồ sơ nhân sự. Cột \oo{Cách tính thuế
TNCN} khi \oo{Nhập Excel} hồ sơ cũng đổi được ô này cho nhiều người một lần.

\subsection{Nội quy công ty}\label{nhan-su-noi-quy}

\duong{Nhân sự \& Lương > Nội quy công ty}, bấm \oo{Thêm tài liệu} để tải file nội quy lên. Mọi người đều đọc
được. Muốn thêm, xoá tài liệu thì vai phải có \oo{Thao tác} ở dòng \oo{Nội quy công ty}.

% ============================================================
\Needspace*{14\baselineskip}
\section{Danh mục dùng chung}\label{danh-muc}

\begin{boicanh}
Năm danh mục đầu (mục \ref{danh-muc-don-vi} đến \ref{danh-muc-cong-doan}) là nguyên liệu của tính giá và lệnh
sản xuất: phiếu tính giá lấy giá giấy, vật tư, công đoạn từ đây; lệnh sản xuất lấy máy, tốc độ, định mức vật tư từ
đây. Khai theo đúng thứ tự: Đơn vị trước, Công đoạn sau cùng, vì công đoạn trỏ tới máy, tổ và vật tư. Sáu danh mục
sau (mục \ref{danh-muc-khuon} đến \ref{danh-muc-muc-thu-chi}) phục vụ kho, KCS, giao hàng, thu chi; khai trước khi
bắt đầu dùng phân hệ đó.
\end{boicanh}

Màn nào cũng cùng một khuôn: thanh trên có nút \oo{Thêm \ldots}, \oo{Nhập Excel}, \oo{Xuất Excel}. Bấm một dòng
trong bảng để mở ngăn sửa, sửa xong bấm \oo{Lưu thay đổi}.

\Needspace*{20\baselineskip}
\subsection{Quyền của các màn danh mục}\label{danh-muc-quyen}

\rong[0.9\linewidth]{40-quyen-danh-muc.png}

\so{1} Mỗi danh mục có \oo{Xem} và \oo{Thao tác} (thêm, sửa, ngừng dùng). \oo{Đơn vị \& quy đổi} chỉ có hai ô này.

\so{2} \so{3} Bốn màn còn lại có quyền chi tiết:
\begin{itemize}
\item \oo{Nhân bản} (Thiết bị \& Máy móc, Công đoạn, Giấy, Vật tư khác): tạo mục mới bằng cách chép một mục có
sẵn. Bản chép mang tên \q{\ldots{} (bản sao)}, nhớ soát lại giá và công thức trước khi dùng.
\item \oo{Xem giá} (Giấy, Công đoạn): giá mua giấy là giá vốn, tiền khoán công đoạn là tiền lương. Không bật thì
các ô đó để trống, người đó sửa dòng vẫn giữ nguyên giá cũ. Riêng giá giấy: ai có \oo{Xem chi tiết giá vốn} của
Tính giá (mục \ref{kinh-doanh-quyen}) vẫn thấy cột \oo{Đơn giá}, vì thẻ sản phẩm ở phiếu tính giá hiện giá từng loại
giấy. Lịch sử giá thì chỉ \oo{Xem giá} của Giấy mới mở được.
\end{itemize}

Sáu màn ở cuối nhóm (\oo{Thành phẩm}, \oo{Khuôn}, \oo{Khai báo kho}, \oo{Tiêu chí KCS}, \oo{Xe giao hàng},
\oo{Mục thu chi}) chỉ có \oo{Xem} và \oo{Thao tác}, không có quyền chi tiết.

\Needspace*{12\baselineskip}
\subsection{Đơn vị \& quy đổi}\label{danh-muc-don-vi}

\begin{boicanh}
Đơn vị khai một lần ở đây, mọi nơi khác (kho, nhà cung cấp, khoán, tính giá) chọn lại từ danh sách này. Quy đổi
cho phần mềm tự đổi giữa các đơn vị, ví dụ m\textsuperscript{2} ra cm\textsuperscript{2}, kg ra tờ.
\end{boicanh}

\canh{0.55}{41-don-vi-them.png}{%
\duong{Cấu hình danh mục > Đơn vị \& quy đổi}, bấm \oo{Thêm đơn vị \& quy đổi}.

\so{1} \oo{Mã}: viết ngắn, không dấu, ví dụ \q{M2}.

\so{2} \oo{Tên}: chữ hiện ra ở mọi màn, ví dụ \q{m\textsuperscript{2}}.

\so{3} Tab \oo{Quy đổi \& số lượng}, khai luôn nếu cần (ảnh dưới).

\so{4} Bấm \oo{Tạo mới}.}

\canh{0.55}{42-don-vi-quy-doi.png}{%
Ví dụ đơn vị cm\textsuperscript{2}. Ở mục \oo{Thêm quy đổi mới}:

\so{1} Gõ số: 1 cm\textsuperscript{2} bằng bao nhiêu đơn vị kia.

\so{2} Chọn đơn vị quy đổi về.

\so{3} Bấm \oo{Thêm quy đổi}.

Quy đổi khai ở một đơn vị thì đơn vị kia tự thấy. Trong ảnh, quy đổi m\textsuperscript{2} và
cm\textsuperscript{2} khai ở m\textsuperscript{2}, nên cm\textsuperscript{2} hiện ở dòng \oo{Quy đổi từ đơn vị
khác}, không cần khai lại.}

Hai đơn vị \oo{kg} và \oo{tờ} là đơn vị bắt buộc của giấy: không xoá, không ngừng dùng, không đổi mã được, quy đổi
của chúng vẫn sửa được. Ở cột \oo{Quy đổi} của bảng, chip màu cam là quy đổi theo công thức (ví dụ tờ ra kg phụ
thuộc định lượng và khổ giấy của từng lệnh), chip xám là quy đổi theo một hệ số cố định.

\Needspace*{22\baselineskip}
\subsection{Thiết bị \& Máy móc}\label{danh-muc-may}

\begin{boicanh}
Máy gắn với tổ đứng máy. Tốc độ và thời gian chuẩn bị khai ở đây là gốc để lệnh sản xuất và xếp lịch tính ra
một bước chạy mất bao lâu.
\end{boicanh}

\canh{0.44}{43-may-chung.png}{%
Bấm \oo{Thêm thiết bị \& máy móc}, hoặc mở một máy có sẵn. Tab \oo{Thông tin chung}:

\so{1} \oo{Nhóm máy}: chọn nhóm, ví dụ Bế, In. Thiếu nhóm thì bấm \oo{Thêm / xoá} cạnh đó.

\so{2} \oo{Tổ đứng máy}: tick các tổ chạy được máy này. Ở danh mục Công đoạn, máy sẽ hiện theo tổ đã tick ở đây.

\oo{Hãng sản xuất}, \oo{Model}, \oo{Số seri} điền nếu có.}

\canh{0.44}{44-may-toc-do.png}{%
Tab \oo{Thông số kỹ thuật}:

\so{1} \oo{Tốc độ trung bình}: số dùng để tính thời gian chạy. \oo{Tốc độ tối thiểu} và \oo{Tốc độ tối đa} cho
lệnh sản xuất thấy khoảng nhanh, chậm. Để trống thì lấy bằng tốc độ trung bình.

\so{2} \oo{Đơn vị tốc độ}: tờ/h hoặc kẽm/h.

\so{3} \oo{Thời gian chuẩn bị}: chọn \oo{Điền tổng} để gõ một số, hoặc \oo{Theo từng khoản} để khai từng việc
và phần mềm tự cộng.

\so{4} Số phút chuẩn bị khi chọn \oo{Điền tổng}.}

\canh{0.44}{68-bao-tri.png}{%
Tab \oo{Lịch bảo trì}: khai các gói bảo trì định kỳ của máy. Mỗi gói có tên (ví dụ Bảo trì tháng máy bế), chu kỳ
\oo{Mỗi} N ngày, tuần, tháng hoặc năm, ngày \oo{Bắt đầu từ}, và danh sách \oo{Việc phải làm trong gói}.

\oo{Kỳ tới} phần mềm tự tính từ phiếu bảo trì gần nhất của máy. Tới kỳ thì phần mềm báo \oo{Tới hạn bảo trì}, quá kỳ mà
chưa có phiếu thì báo \oo{Quá hạn bảo trì}.

Một máy có nhiều gói thì bấm \oo{Thêm gói bảo trì}, ví dụ gói tháng và gói năm.}

\Needspace*{22\baselineskip}
\subsection{Giấy}\label{danh-muc-giay}

\begin{boicanh}
Mỗi dòng là một loại giấy theo định lượng. Đơn giá và công thức ở đây là chỗ phiếu tính giá lấy ra tiền giấy.
\end{boicanh}

\canh{0.5}{45-giay.png}{%
Bấm \oo{Thêm giấy}, hoặc mở một dòng có sẵn.

\so{1} \oo{Loại}: \oo{Giấy} hoặc \oo{Sóng}. Một sản phẩm tính giá, một lệnh có tối đa một Giấy và một Sóng.

\so{2} \oo{Định lượng}, tính bằng g/m\textsuperscript{2}.

\so{3} \oo{Đơn giá}: giá một kg. Khi tính giá, phần mềm tự quy về giá một tờ hoặc một kg.

Giấy không có ô chọn đơn vị: phần mềm tự gán hai đơn vị kg và tờ, đổi qua lại theo khổ và
định lượng.

\so{4} Tab \oo{Công thức tính giá} (ảnh dưới).}

\canh{0.5}{46-giay-cong-thuc.png}{%
\so{1} Ô công thức. Dòng \oo{Chèn toán tử} có phép tính và hàm như \oo{round}, \oo{max}, \oo{if}. Chế độ
\oo{Trực quan} hiện công thức bằng thẻ, \oo{Dạng chữ} để gõ tay.

\so{2} \oo{Danh sách biến khả dụng}: bấm một biến là chèn vào công thức. Ví dụ trong ảnh: định lượng × dài × rộng
tờ nguyên × số tờ × đơn giá.}

\Needspace*{22\baselineskip}
\subsection{Vật tư khác}\label{danh-muc-vat-tu}

\begin{boicanh}
Mọi vật tư không phải giấy: mực, keo, màng, kẽm, dụng cụ. Mỗi vật tư có hai công thức: một công thức ra tiền ở
phiếu tính giá, một công thức ra lượng tiêu hao để lệnh sản xuất xin vật tư.
\end{boicanh}

\canh{0.5}{47-vat-tu.png}{%
Bấm \oo{Thêm vật tư khác}, hoặc mở một dòng có sẵn.

\so{1} \oo{Đơn vị tính}: đơn vị chính. Đơn giá của vật tư tính theo đơn vị này.

\so{2} \oo{Đơn vị thêm}: đơn vị kho hoặc nhà cung cấp còn dùng, ví dụ thùng, cuộn. Giá theo đơn vị thêm lấy ở
bảng giá nhà cung cấp.

\so{3} \oo{Chip riêng của vật tư}: con số người tính giá phải nhập cho riêng vật tư này. Ví dụ \q{Diện tích phủ UV
(\%)} thành biến \q{ty\_le\_phu} dùng trong hai công thức.

\so{4} \oo{Công thức tính giá}: ra tiền của vật tư ở phiếu tính giá.

\so{5} \oo{Công thức định mức}: ra lượng vật tư tiêu hao cho một bước. Lệnh sản xuất tự tính, không nhập tay.}

\canh{0.5}{48-vat-tu-dinh-muc.png}{%
Hai tab công thức dùng chung cách khai với Giấy:

\so{1} Ô công thức.

\so{2} Danh sách biến. Chip riêng của vật tư nằm cuối danh sách, ví dụ \oo{Diện tích phủ UV (\%)}.

Vật tư chưa có công thức thì bảng ghi \oo{Chưa có tính giá} hoặc \oo{Chưa có định mức}.}

\Needspace*{22\baselineskip}
\subsection{Công đoạn}\label{danh-muc-cong-doan}

\begin{boicanh}
Công đoạn là một bước làm ra sản phẩm: in, cán màng, bế, dán. Nó nối tổ, máy, vật tư và cách tính tiền lại với
nhau, nên khai sau cùng. Tab \oo{Khoán} là tiền công trả thợ, viết riêng ở mục \ref{danh-muc-khoan}.
\end{boicanh}

\canh{0.5}{49-cong-doan.png}{%
Bấm \oo{Thêm công đoạn}, hoặc mở một dòng có sẵn. Tab \oo{Thông tin}:

\so{1} \oo{Giai đoạn}, ví dụ Gia công sau in.

\so{2} \oo{Phòng ban / Tổ phụ trách}: tick được nhiều tổ. Lệnh sản xuất chọn một trong số này cho từng bước. Tổ
tick đầu tiên là tổ mặc định.}

\canh{0.5}{50-cong-doan-may.png}{%
Kéo xuống:

\so{1} \oo{Bước này cần khuôn}: bật nếu bước này dùng khuôn.

\so{2} \oo{Loại khuôn}, ví dụ khuôn bế.

\so{3} \oo{Máy chạy được công đoạn này}: máy hiện theo tổ đã tick ở trên. Tốc độ chỉ đọc, muốn sửa thì sửa ở
danh mục máy (mục \ref{danh-muc-may}). Bấm \oo{Đặt mặc định} dưới máy hay dùng nhất.}

\canh{0.5}{53-cong-doan-don-vi.png}{%
Cuối tab:

\so{1} \oo{Đơn vị đầu vào}: thứ đưa vào bước, ví dụ Tờ in.

\so{2} \oo{Đơn vị đầu ra}: thứ bước làm ra, ví dụ Con.

\so{3} \oo{Bù hao}: \oo{Không bù hao}, \oo{Theo bậc số lượng}, hoặc \oo{Cộng cố định (số tờ)}.}

\rong[0.8\linewidth]{67-bu-hao-bac.png}

Bù hao là số tờ hỏng phải chạy thêm ở bước này. \so{1} Chọn \oo{Theo bậc số lượng} thì hiện \so{2} bảng bậc: mỗi
dòng là \q{khi số lượng ra của công đoạn đến N thì bù hao cộng thêm X}, X tính bằng tờ hoặc \% (bấm chữ \oo{Tờ}
để đổi). Dòng cuối \oo{Trên} áp cho mọi số lớn hơn. Chọn \oo{Cộng cố định (số tờ)} thì chỉ gõ một số tờ cộng thêm
cho mọi lệnh. Phần mềm cộng bù hao ngược từ bước cuối về bước in, nên số tờ vào máy in đã gồm bù hao của mọi bước
sau.

\canh{0.5}{51-cong-doan-vat-tu.png}{%
Tab \oo{Vật tư}:

\so{1} Gõ tên vật tư bước này tiêu thụ rồi chọn. Lượng tiêu hao không khai ở đây mà lấy từ công thức định mức của
vật tư (mục \ref{danh-muc-vat-tu}).}

\canh{0.5}{52-cong-doan-cong-thuc.png}{%
Tab \oo{Công thức tính giá}: ra tiền của bước này ở phiếu tính giá, khai giống Giấy.

\so{1} Ô công thức.

\so{2} Danh sách biến. Công đoạn có thêm \oo{SL vào của công đoạn} và \oo{SL ra của công đoạn}. Nghĩa từng
biến ở mục \ref{phu-luc-chip}.}

\Needspace*{20\baselineskip}
\subsection{Tiền khoán của công đoạn}\label{danh-muc-khoan}

\begin{boicanh}
Tab \oo{Khoán} quyết định mỗi lần thợ ghi sản lượng ở bước này được bao nhiêu tiền. Tiền này cộng vào lương khoán
(mục \ref{nhan-su-cau-hinh-luong}). Tab chỉ hiện số tiền cho người có ô \oo{Xem giá} của Công đoạn.
\end{boicanh}

\rong[0.85\linewidth]{64-khoan-in.png}

\so{1} \oo{Đơn giá}: \oo{Đơn vị tính khoán}, \oo{Giá cố định chung}, \oo{Đơn giá chung}. Đây là giá chung cho mọi
tổ, trừ tổ có giá riêng. \so{2} \oo{Bậc giá theo số màu}: dùng cho bước in. Mỗi dòng có khoảng \oo{Từ màu}
\oo{Đến màu}, điều kiện \oo{Chỉ khi có pha} hoặc \oo{Chỉ khi tick}, và giá cố định, đơn giá riêng từng tổ. Lúc
ghi sản lượng, phần mềm dò từ trên xuống, lệnh hợp dòng nào trước thì ăn giá dòng đó. Số màu là số kẽm của lệnh,
như bên Tính giá. Ô giá trống là tổ không chạy bậc này.

\rong[0.85\linewidth]{65-khoan-to.png}

\so{1} \oo{Giá theo tổ}: để trống hết thì mọi tổ lấy giá chung. Đã điền cho một tổ thì tổ còn trống bị coi là
thiếu giá. \so{2} \oo{Tuỳ chọn khi ghi}: thứ thợ chọn mỗi lần ghi, ví dụ số lỗ, kẽm bế nối. Mỗi dòng có
\oo{Tên hiện cho thợ}, cách nhập (\oo{Tick ô} thì tick là 1, không tick là 0; \oo{Gõ số} thì bỏ trống là 0) và
gợi ý cho thợ. Mỗi tuỳ chọn thành một biến dùng trong công thức khoán. \so{3} \oo{Việc phát sinh}: việc lặt vặt
trả tiền riêng, không cộng vào sản lượng. Khai tên việc, đơn giá, đơn vị tính, cách tính giá theo từng tổ.

\rong[0.85\linewidth]{66-khoan-cong-thuc.png}

\so{1} \oo{Công thức khoán}: tính trên từng lần ghi sản lượng, giá lấy theo tổ của lần ghi. Kết quả công thức là
tiền của lần ghi, phần mềm không nhân thêm gì. Bỏ trống thì lần ghi được 0 đ, phần mềm không tự lấy đơn giá nhân
sản lượng. Công thức dùng được biến khoán (mục \ref{phu-luc-chip}) và các tuỳ chọn khi ghi. Ví dụ bước bế khuôn
đứt có thể viết \texttt{sl\_ghi / max(1, so\_con) * don\_gia\_khoan}: số tờ bế nhân đơn giá khoán.

\Needspace*{20\baselineskip}
\subsection{Khuôn}\label{danh-muc-khuon}

\begin{boicanh}
Kho khuôn bế của xưởng. Mỗi khuôn làm riêng cho hình của một sản phẩm, đơn lặp lại thì lấy khuôn cũ ra dùng. Khai
để tìm lại được: của khách nào, cất ở kệ nào, còn dùng được không. Phiếu tính giá không đọc danh mục này, chỉ hỏi
khuôn có sẵn hay làm mới.
\end{boicanh}

\canh{0.45}{57-khuon.png}{%
\duong{Cấu hình danh mục > Khuôn}, bấm \oo{Thêm khuôn}.

\so{1} \oo{Tên}: tên sản phẩm của khuôn, ví dụ Khuôn bế hộp mỹ phẩm 8×8×12.

\so{2} \oo{Khách hàng}: khuôn của khách nào.

\so{3} \oo{Loại}: hiện chỉ có \oo{Khuôn bế}.

\so{4} \oo{Số kệ / vị trí lưu}.

\so{5} \oo{Tình trạng}: \oo{Đang dùng}, \oo{Đang đặt làm}, \oo{Hỏng}, \oo{Thanh lý}.

Ở lệnh sản xuất, bước có bật \oo{Bước này cần khuôn} (mục \ref{danh-muc-cong-doan}) chọn khuôn từ đây, lọc theo
khách và loại. Chọn làm khuôn mới thì phần mềm tự thêm một dòng ở đây với tình trạng \oo{Đang đặt làm}.}

\Needspace*{16\baselineskip}
\subsection{Thành phẩm}\label{danh-muc-thanh-pham}

\begin{boicanh}
Thành phẩm là hàng của đơn hàng bán. Phần mềm tự khai khi chốt đơn, không có nút thêm tay. Màn này để xem, sửa
tên và đơn vị tính.
\end{boicanh}

\rong[0.9\linewidth]{55-thanh-pham.png}

\so{1} \oo{Đơn đầu tiên}: đơn hàng bán đã sinh ra thành phẩm này. \so{2} \oo{Nhập Excel}: chỉ sửa hàng loạt tên,
đơn vị tính của mã đã có. Dòng mang mã chưa có sẽ báo lỗi, vì thành phẩm mới chỉ sinh ra từ đơn hàng bán.

\canh{0.45}{56-thanh-pham-sua.png}{%
Bấm một dòng để mở:

\so{1} \oo{Tên}, sửa được.

\so{2} \oo{Đơn vị tính}: cái, hộp, thùng.

\so{3} \oo{Đơn hàng đầu tiên} và \oo{Khách hàng}: chỉ đọc.

Thành phẩm thuộc về một khách. Khách khác đặt cùng tên thì phần mềm tạo một mã khác.}

\Needspace*{16\baselineskip}
\subsection{Khai báo kho}\label{danh-muc-kho}

\begin{boicanh}
Mỗi kho khai ở đây tự hiện thành một mục dưới nhóm \oo{Kho hàng} trên thanh bên. \oo{Kho thành phẩm} (mã
\oo{KHO-TP}) là kho cố định của phần mềm: thành phẩm chỉ nhập và xuất ở kho này.
\end{boicanh}

\rong[0.8\linewidth]{58-kho.png}

\duong{Cấu hình danh mục > Khai báo kho}, bấm \oo{Thêm kho hàng}. \so{1} \oo{Tên}, ví dụ Kho giấy. \so{2}
\oo{Vị trí kho}: kho nằm ở đâu, ví dụ Xưởng 2. Bấm \oo{Tạo mới}.

\so{3} Tab \oo{Vị trí kho}: khai kệ, ô cất hàng trong kho (gõ tên như Kệ A1, bấm \oo{Thêm}). Lập lô hoặc phiếu
nhập xuất sẽ chọn vị trí từ danh sách này. Tab này chỉ dùng được sau khi kho đã lưu: mở lại kho vừa tạo rồi khai.

\Needspace*{20\baselineskip}
\subsection{Tiêu chí KCS}\label{danh-muc-kcs}

\begin{boicanh}
Tiêu chí KCS là các điểm phải kiểm sau mỗi công đoạn, ví dụ màu đúng mẫu, bế không xơ mép. Lệnh sản xuất lúc
phát hành chép bộ tiêu chí của từng công đoạn về để người kiểm đánh dấu.
\end{boicanh}

\rong[0.78\linewidth]{60-kcs.png}

\so{1} Cột trái là cây công đoạn theo giai đoạn. Công đoạn chưa có tiêu chí ghi \oo{chưa có}. Bấm một công đoạn
để khai. \so{2} Gõ một tiêu chí rồi bấm Enter, lặp lại cho tiêu chí tiếp theo. \so{3} \oo{Chép từ công đoạn
khác}: lấy nguyên bộ tiêu chí của một công đoạn giống. \so{4} \oo{Xem KCS cuối thấy gì}: chọn các công đoạn một
lệnh sẽ chạy qua, xem bước KCS cuối sẽ bày những tiêu chí nào.

Sửa tiêu chí chỉ áp cho lệnh phát hành từ lúc đó. Lệnh đã phát hành giữ bộ cũ, hồ sơ lệnh sẽ báo tiêu chí đã đổi
và có nút \oo{Cập nhật tiêu chí}.

\Needspace*{16\baselineskip}
\subsection{Xe giao hàng}\label{danh-muc-xe}

\canh{0.45}{61-xe.png}{%
Khai xe để tính tiền khoán km cho tài xế. Cần khai mức khoán km ở Cấu hình lương trước (mục
\ref{nhan-su-cau-hinh-luong}).

\duong{Cấu hình danh mục > Xe giao hàng}, bấm \oo{Thêm xe giao hàng}.

\so{1} \oo{Biển số}.

\so{2} \oo{Tên}, ví dụ Xe anh Tùng.

\so{3} \oo{Mức khoán km}: chọn mức đã khai. Giá mỗi km do mức này quyết định.

\so{4} \oo{Tải trọng}: chỉ để đối chiếu, không ảnh hưởng tiền.}

\Needspace*{20\baselineskip}
\subsection{Mục thu chi}\label{danh-muc-muc-thu-chi}

\begin{boicanh}
Mỗi phiếu thu, phiếu chi mang một mục. Báo cáo thu chi gom tiền theo các mục này.
\end{boicanh}

\rong[0.8\linewidth]{62-muc-list.png}

\so{1} Các dòng mã \oo{HT-\ldots} có ổ khoá là mục của phần mềm, ví dụ Trả nhà cung cấp, Khách trả tiền hàng.
Phần mềm tự gắn khi lập phiếu từ đơn mua, đơn bán; không sửa, không xoá được. \so{2} Cột \oo{Bên}: \oo{Thu} hoặc
\oo{Chi}. Cột \oo{Đã dùng} đếm số phiếu đang mang mục đó.

\canh{0.45}{63-muc.png}{%
Thêm mục riêng của công ty, ví dụ Tiền điện, Bán phế liệu: bấm \oo{Thêm mục thu chi}.

\so{1} \oo{Mã}: phần mềm gợi ý sẵn, sửa được.

\so{2} \oo{Tên}.

\so{3} \oo{Loại}: \oo{Thu} hoặc \oo{Chi}.

Mục tự thêm mà chưa phiếu nào dùng thì xoá được.}

% ============================================================
\Needspace*{14\baselineskip}
\section{Kinh doanh}\label{kinh-doanh}

\begin{boicanh}
Giá trên phiếu tính giá lấy từ danh mục ở chương \ref{danh-muc}. Phần cấu hình còn lại của phòng Kinh doanh gồm
ba việc: cấp quyền cho từng màn, khai khách hàng kèm chính sách tài chính, rồi lập một phiếu tính giá mẫu để kiểm
cả chuỗi. Cách lập phiếu, báo giá, đơn hàng xem ở tài liệu hướng dẫn sử dụng phòng Kinh doanh.
\end{boicanh}

\Needspace*{16\baselineskip}
\subsection{Quyền của sáu màn}\label{kinh-doanh-quyen}

Nhóm \duong{Kinh doanh} có sáu màn. \oo{Quy trình kinh doanh} là sơ đồ luồng bán hàng, chỉ cần \oo{Xem}. Năm màn
còn lại có quyền chi tiết. Ảnh dưới chụp vai Trợ lý kinh doanh: \so{1} bấm chip \oo{chi tiết} để mở, \so{2} các
quyền chi tiết của dòng.

\rong[0.7\linewidth]{70-quyen-khach.png}

\oo{Khách hàng}: \oo{Thao tác} là thêm, sửa khách, nhập Excel, người liên hệ, lịch hẹn. Phần mềm không có xoá
khách. \oo{Điều chuyển} cho đổi người phụ trách (mục \ref{kinh-doanh-dieu-chuyen}). \oo{Thiết lập chính sách
tài chính} cho sửa hạn mức công nợ và rào giá của khách (mục \ref{kinh-doanh-tai-chinh}); không bật thì chỉ xem.

\rong[0.7\linewidth]{71-quyen-tinh-gia.png}

\oo{Tính giá}: \oo{Xem chi tiết giá vốn} cho thấy từng dòng giá vốn và thẻ sản phẩm (giấy, khổ, công đoạn). Không
bật thì chỉ thấy giá vốn tổng. Có \oo{Thao tác} thì ô này tự bật. Sale thường chỉ để \oo{Xem}.

\rong[0.7\linewidth]{72-quyen-bao-gia.png}

\oo{Báo giá in ấn}: \oo{Thao tác} gồm tạo, sửa, gửi khách, ghi kết quả, huỷ, tải PDF. \oo{Duyệt báo giá đặc thù}
dành cho người duyệt (thường là giám đốc). Báo giá đặc thù là báo giá từ 1 tỷ trở lên (trước VAT), bán dưới giá
vốn, hoặc chiết khấu, markup, kê ngoài rào của khách. \oo{Tạo đơn hàng} là ô chung với \oo{Thao tác} của dòng
Đơn hàng bán: bật một chỗ thì chỗ kia bật theo.

\rong[0.7\linewidth]{73-quyen-don-ban.png}

\oo{Đơn hàng bán}: \oo{Thao tác} gồm tạo, sửa, chốt, huỷ đơn và chuyển đơn xuống sản xuất. Ba quyền chi tiết
dính tới tiền, nên chung ô với màn kế toán: \oo{Thu cọc và ghi hoá đơn} chung với Phiếu thu, \oo{Huỷ hoá đơn}
chung với Hủy phiếu thu, \oo{Trả lại cọc} chung với Phiếu chi.

\rong[0.7\linewidth]{74-quyen-giao-hang.png}

\oo{Giao hàng}: \oo{Thao tác} là gửi yêu cầu giao, báo đã lấy hàng, nhập kết quả, số km, đính biên bản. \oo{Lên
đơn giao hàng} cho xếp chuyến và gửi đề nghị xuất sang kho. \oo{Huỷ yêu cầu / huỷ chuyến} phải ghi lý do.
\oo{Nhân viên giao hàng} cho xem lịch, chuyến, số km của tài xế khác. Phạm vi ở dòng này tính riêng: \oo{Của tôi}
là yêu cầu mình lập hoặc chuyến mình chở. Kế toán cần đính biên bản cho mọi đơn thì để \oo{Tất cả}.

Phạm vi các dòng khác theo luật chung (mục \ref{phan-quyen-pham-vi}). Sale để \oo{Của tôi} và lập nhóm dùng chung
với trợ lý của mình (mục \ref{phan-quyen-nhom}) thì hai người cùng thấy khách, báo giá, đơn của nhau.

\Needspace*{20\baselineskip}
\subsection{Khai khách hàng}\label{kinh-doanh-khach}

\begin{boicanh}
Sale khai khách của mình. Người phụ trách khách là người hưởng hoa hồng khi đơn của khách thu được tiền, kể cả
khi trợ lý lập phiếu, báo giá, đơn hộ.
\end{boicanh}

\canh{0.55}{75-kh-them.png}{%
\duong{Kinh doanh > Khách hàng}, bấm \oo{Tạo khách hàng}.

\so{1} Tên công ty hoặc tên khách, bắt buộc.

\so{2} Gõ mã số thuế rồi bấm \oo{Tra cứu MST}: phần mềm tự điền tên và địa chỉ.

\so{3} \oo{Công ty} hoặc \oo{Cá nhân}.

\so{4} \oo{NV phụ trách}: để trống là chính người tạo.

\so{5} \oo{Tạo khách hàng}.}

Trùng mã số thuế, tên hoặc email với khách đã có thì phần mềm cảnh báo nhưng vẫn cho lưu. Đọc kỹ trước khi bấm
tiếp: khách trùng là hai hồ sơ riêng, báo giá và công nợ bị tách ra hai nơi. Khách đang có trong sổ sách cũ thì
dùng \oo{Nhập Excel} để đưa vào một lần.

\Needspace*{12\baselineskip}
\subsection{Chính sách tài chính của khách}\label{kinh-doanh-tai-chinh}

\begin{boicanh}
Khách mới tạo chưa có hạn mức và rào giá. Người có quyền \oo{Thiết lập chính sách tài chính} (thường là trưởng
phòng hoặc kế toán) đặt cho từng khách. Rào giá quyết định báo giá nào phải trình duyệt.
\end{boicanh}

\rong[0.7\linewidth]{77-kh-taichinh.png}

Mở hồ sơ khách, tab \oo{Tổng quan}, kéo xuống cuối: \so{1} thẻ \oo{Chính sách tài chính}. \so{2} Bấm \oo{Sửa}.

\canh{0.55}{78-kh-taichinh-sua.png}{%
\so{1} \oo{Hạn mức công nợ}: Công nợ phải thu đánh dấu khách nợ vượt số này.

\so{2} \oo{Hạn thanh toán}: số ngày kể từ ngày xuất hoá đơn. Quá ngày thì khoản nợ tính là quá hạn.

\so{3} \oo{Rào giá khi báo giá}: khoảng \oo{Từ} đến \oo{Đến} của \oo{Chiết khấu cho phép}, \oo{Markup trên giá
vốn}, \oo{Kê cho phép}. Báo giá ra ngoài khoảng thì phải trình duyệt. Để trống một đầu là không giới hạn đầu đó.

\oo{Thuế TNDN trên kê}: phần trăm giữ lại để nộp thuế, chi khách phần còn lại.

Bấm \oo{Lưu}.}

\Needspace*{10\baselineskip}
\subsection{Đổi người phụ trách}\label{kinh-doanh-dieu-chuyen}

\canh{0.5}{76-kh-dieu-chuyen.png}{%
Khi sale nghỉ hoặc chia lại khách: \duong{Kinh doanh > Khách hàng}, bấm \oo{Điều chuyển}. Cần quyền
\oo{Điều chuyển}.

\so{1} Nhân viên đang phụ trách.

\so{2} Nhân viên nhận.

Toàn bộ khách của người nguồn chuyển sang người nhận, kèm lịch sử giao dịch. Nút \oo{Xác nhận bàn giao} đếm
ngược vài giây trước khi bấm được.}

\Needspace*{6\baselineskip}
\subsection{Kiểm bằng một phiếu tính giá mẫu}\label{kinh-doanh-kiem}

Khai xong danh mục và khách, trợ lý kinh doanh lập một phiếu tính giá cho một sản phẩm đã từng báo giá tay, đủ
các bước quen thuộc, ví dụ hộp giấy in 4 màu, cán màng, bế, dán. So từng dòng chi phí trên phiếu với bảng tính
tay:

\begin{itemize}
\item Tiền giấy lệch: kiểm \oo{Đơn giá} và \oo{Định lượng} của giấy (mục \ref{danh-muc-giay}), khổ giấy chọn trên
phiếu.
\item Một công đoạn ra 0 đ: công đoạn đó chưa có công thức tính giá (mục \ref{danh-muc-cong-doan}).
\item Số tờ in lệch: kiểm \oo{Đơn vị đầu vào}, \oo{Đơn vị đầu ra} và \oo{Bù hao} của từng công đoạn trong chuỗi.
\item Tiền vật tư thiếu: vật tư chưa chọn ở tab \oo{Vật tư} của công đoạn, hoặc vật tư chưa có công thức tính giá
(mục \ref{danh-muc-vat-tu}).
\end{itemize}

Mọi dòng khớp với cách tính tay thì phần cấu hình ban đầu xong.

\clearpage
\section{Bảng chip trong công thức}\label{phu-luc-chip}

Ô công thức nào cũng có danh sách chip bên dưới, bấm chip để chèn. Mỗi chip là một con số phần mềm tự điền lúc
tính. Cột cuối cho biết chip có mặt ở ô nào: G là công thức giá của Giấy, V là Vật tư khác, C là công thức tính
giá của Công đoạn, Q là công thức lượng (định mức vật tư, quy đổi đơn vị), K là công thức khoán. Chip đặt sai ô
thì phần mềm không cho lưu.

{\small
\begin{longtable}{@{}p{0.27\linewidth}p{0.6\linewidth}l@{}}
\toprule
\tieude Chip & \tieude Nghĩa & \tieude Ô\\
\endhead
\multicolumn{3}{l}{\tieude\small Kích thước và số lượng}\\
\oo{Dài / Rộng sản phẩm} & Khổ thành phẩm, tính bằng mét (gõ 210 mm thì phần mềm đổi ra 0,21). & GVCQK\\
\oo{Dài / Rộng tờ nguyên} & Khổ tờ giấy to chưa cắt, mét. & GVCQK\\
\oo{Dài / Rộng tờ in} & Khổ tờ chạy máy in, mét. & GVCQK\\
\oo{Số lượng đặt} & Số cái khách đặt. & GVCQK\\
\oo{Số con/tờ in} & Một tờ in ra mấy cái. & GVCQK\\
\oo{Số trang sách} & Số trang nội dung một sản phẩm. Tờ rời, hộp là 1. & GVCQK\\
\oo{Trang mỗi tay} & Số trang trên một tay gấp (8, 16, 32). Không gấp tay là 1. & GVCQK\\
\midrule
\multicolumn{3}{l}{\tieude\small Số tờ}\\
\oo{Tờ vào máy} & Tờ cần in cộng bù hao, đã gồm bù hao các bước sau. & GVCQK\\
\oo{Tờ tốt sau in} & Số tờ tốt sau in, dùng cho gia công. & GVCQK\\
\oo{Tờ giấy nguyên} & Số tờ giấy to tiêu hao, chưa cắt. & GVCQK\\
\midrule
\multicolumn{3}{l}{\tieude\small Màu, mặt, kẽm}\\
\oo{Số màu in} & Số màu lên máy, bằng số kẽm một tay, đổi theo quy cách in. & GVCQK\\
\oo{Số màu pha} & Số kẽm màu pha (Pantone) nằm trong Số màu in. & GVCQK\\
\oo{Số màu phủ giấy} & Tổng màu phủ lên giấy, mặt A cộng mặt B. Dùng cho mực, giờ máy. & GVCQK\\
\oo{Số mặt in} & 1 mặt là 1; 2 mặt hoặc tự trở là 2. & GVCQK\\
\oo{Số bản kẽm} & Số bản mỗi tay nhân số tay. & GVCQK\\
\midrule
\multicolumn{3}{l}{\tieude\small Giấy, giá, số của bước}\\
\oo{Định lượng giấy} & kg/m\textsuperscript{2}, bằng gsm chia 1.000. & GVQK\\
\oo{Đơn giá giấy} & Giá một kg của loại giấy đang tính. & G\\
\oo{SL vào / ra của công đoạn} & Số lượng vào, ra của chính bước đang tính, theo đơn vị đầu vào, đầu ra của bước. & VCQK\\
\oo{Số lượt qua máy} & Số lần hàng đi qua bước này. In trở 2 mặt là 2. & VQK\\
\oo{Đơn giá khoán} & Đơn giá của chính đầu việc khoán đang khai. Dùng chip này thì công thức ra tiền. & QK\\
\midrule
\multicolumn{3}{l}{\tieude\small Chỉ có ở công thức khoán}\\
\oo{SL ghi} & Số lượng thợ ghi ở lần ghi. & K\\
\oo{Giá cố định} & Giá cố định của tổ ghi, ví dụ giá một bài in. & K\\
\oo{Số lần ghi sản lượng} & Số lần ghi của cùng lệnh ở cùng công đoạn tới cuối kỳ lương, tối thiểu 1. & K\\
\oo{Số lần ghi có thử việc} & Trong số lần ghi trên, số lần có người đang thử việc ngày ghi. & K\\
\oo{Chạy tiếp} & 1 khi lần ghi liền trước trên cùng máy là chính việc này và cùng ngày, ngược lại 0. & K\\
\oo{Trùng khuôn} & 1 khi lần ghi liền trước trên cùng máy dùng cùng con dao, khỏi lên khuôn. & K\\
\oo{Khuôn mới} & 1 khi đây là lần ghi đầu tiên dùng con dao của bước. & K\\
\oo{Tiền giờ} & Tổng (lương vị trí + trách nhiệm) của người làm chia 26, chia 8, nhân số giờ lần ghi. & K\\
\oo{Tự trở} & 1 khi quy cách in là tự trở hoặc trở nhíp. & K\\
\oo{Giá đã gồm UV} & 1 khi bậc giá của lần ghi đã gồm mực UV. & K\\
\bottomrule
\end{longtable}}

Ngoài bảng trên, công thức khoán còn dùng được các \oo{Tuỳ chọn khi ghi} của chính công đoạn (mục
\ref{danh-muc-khoan}). Ô dạng 0/1 hay dùng với \texttt{if}, ví dụ \texttt{if(khuon\_moi > 0, 80000, gia\_co\_dinh)}.

\end{document}
